\documentclass[aspectratio=169, 11pt]{beamer}

% ── Theme & Colors (matched from enzyme_agent_workflow.tex) ──
\usetheme{Madrid}
\usecolortheme{default}

\definecolor{deepblue}{RGB}{20, 66, 114}
\definecolor{accentblue}{RGB}{41, 128, 185}
\definecolor{lightblue}{RGB}{214, 234, 248}
\definecolor{accentorange}{RGB}{230, 126, 34}
\definecolor{accentgreen}{RGB}{39, 174, 96}
\definecolor{darkgray}{RGB}{52, 73, 94}
\definecolor{lightgray}{RGB}{236, 240, 241}
\definecolor{accentred}{RGB}{192, 57, 43}
\definecolor{softpurple}{RGB}{142, 68, 173}

\setbeamercolor{structure}{fg=deepblue}
\setbeamercolor{palette primary}{bg=deepblue, fg=white}
\setbeamercolor{palette secondary}{bg=accentblue, fg=white}
\setbeamercolor{palette tertiary}{bg=deepblue!80, fg=white}
\setbeamercolor{frametitle}{bg=deepblue, fg=white}
\setbeamercolor{title}{fg=white}
\setbeamercolor{subtitle}{fg=lightblue}
\setbeamercolor{block title}{bg=accentblue, fg=white}
\setbeamercolor{block body}{bg=lightblue!30}
\setbeamercolor{block title alerted}{bg=accentorange, fg=white}
\setbeamercolor{block body alerted}{bg=accentorange!10}
\setbeamercolor{block title example}{bg=accentgreen, fg=white}
\setbeamercolor{block body example}{bg=accentgreen!10}
\setbeamercolor{item}{fg=accentblue}

\setbeamertemplate{navigation symbols}{}
\setbeamertemplate{footline}{%
  \hfill\insertframenumber/\inserttotalframenumber\hspace*{4pt}\vspace{2pt}%
}

% ── Packages ────────────────────────────────────────────────
\usepackage{tikz}
\usetikzlibrary{shapes.geometric, arrows.meta, positioning, calc, fit, backgrounds}
\usepackage{booktabs}
\usepackage{graphicx}
\usepackage{colortbl}
\usepackage{adjustbox}
\usepackage{multirow}

\graphicspath{{../experimental_feedback/}}

\title{Structure-Aware Enzyme--Substrate Co-folding}
\subtitle{From Candidate List to Structural Evidence}
\author{R. Zhu}
\date{\today}

\begin{document}

% ════════════════════════════════════════════════════════════
% SLIDE 1: Title
% ════════════════════════════════════════════════════════════
{
\setbeamercolor{background canvas}{bg=deepblue}
\begin{frame}[plain]
\vfill
\begin{center}
{\Huge\bfseries\color{white} Structure-Aware Co-folding}\\[10pt]
{\Large\color{lightblue} Predicting Enzyme--Substrate Complexes\\for Non-Natural Reactions}\\[24pt]
{\large\color{white} R. Zhu}\\[6pt]
{\normalsize\color{lightblue!70} \today}
\end{center}
\vfill
\end{frame}
}

% ════════════════════════════════════════════════════════════
% SLIDE 2: Motivation
% ════════════════════════════════════════════════════════════
\begin{frame}{Why Structural Modeling?}
\begin{columns}[T]
\begin{column}{0.48\textwidth}
\begin{alertblock}{The Problem}
\begin{itemize}
\item Discovery agent ranks by \textbf{sequence \& literature}
\item But: does the substrate \textbf{physically fit}?
\item Is the pocket \textbf{chemically compatible}?
\item Is the substrate \textbf{positioned for catalysis}?
\end{itemize}
\end{alertblock}
\vspace{4pt}
\colorbox{accentred!15}{\parbox{0.92\textwidth}{\centering\small
\textbf{McbA:} iPTM = 0.987, hydrophobicity = 42.9\%\\
\textit{yet experimentally inactive}
}}
\end{column}
\begin{column}{0.48\textwidth}
\begin{exampleblock}{Our Approach}
\begin{itemize}
\item \textbf{Co-fold} each candidate with substrates + cofactors using Boltz-2
\item Extract \textbf{binding geometry}, \textbf{pocket chemistry}, \textbf{pocket stability}
\item Compare \textbf{apo vs.\ holo} pockets
\item Independent metrics --- \textbf{no composite score}
\end{itemize}
\end{exampleblock}
\end{column}
\end{columns}

\vspace{4pt}
\begin{center}
\colorbox{lightblue!50}{\parbox{0.88\textwidth}{\centering
\small\textbf{Goal:} Structural evidence for \textit{which} candidates to test and \textit{why} others should be deprioritized.
}}
\end{center}
\end{frame}

% ════════════════════════════════════════════════════════════
% SLIDE 3: Pipeline overview
% ════════════════════════════════════════════════════════════
\begin{frame}{Co-folding Pipeline Overview}
\vspace{-8pt}
\begin{center}
\includegraphics[height=0.72\textheight]{figures/cofold_pipeline.png}
\end{center}
\vspace{-6pt}
\begin{center}
\scriptsize
\textbf{Input:} 29 synthetase + 93 hydrolase candidates $\rightarrow$
\textbf{Output:} Structurally scored candidates across 8 metrics
\end{center}
\end{frame}

% ════════════════════════════════════════════════════════════
% SLIDE 4: Phase 1+2 — Pocket + Catalytic Residues
% ════════════════════════════════════════════════════════════
\begin{frame}{Phases 1--2: Pocket Detection \& Catalytic Residue Mapping}
\vspace{-4pt}
\begin{columns}[T]
\begin{column}{0.48\textwidth}
\begin{block}{Phase 1: Where Is the Pocket?}
Three complementary methods:
\begin{enumerate}\small
\item \textbf{AlphaFill} --- transplant cofactors from PDB homologs into AlphaFold structures\\
{\scriptsize (filters artifacts: GOL, PEG, buffer salts)}
\item \textbf{pyKVFinder} --- geometric cavity detection
\item \textbf{P2Rank} --- ML-based prediction (fallback)
\end{enumerate}
\vspace{2pt}
Pocket selected by \textbf{overlap with catalytic residues} from Phase 2
\end{block}
\end{column}

\begin{column}{0.48\textwidth}
\begin{block}{Phase 2: Which Residues Are Catalytic?}
Two-phase identification:

\textbf{A: Direct annotation} (fast)
\begin{itemize}\small
\item UniProt ``Active site'' $\rightarrow$ core
\item InterPro homolog consensus ($\geq$50\%)
\end{itemize}

\textbf{B: FoldSeek + LLM} (fills gaps)
\begin{itemize}\small
\item Structural homologs at TM $\geq$ 0.7
\item Transfer positions via C$\alpha$ alignment
\item \textcolor{accentorange}{\textbf{LLM validates}} biochemical plausibility
\end{itemize}
\end{block}

\vspace{2pt}
\colorbox{accentgreen!15}{\parbox{0.92\textwidth}{\centering\small
Phase B raised hydrolase catalytic coverage from \textbf{3\%} to \textbf{24\%}
}}
\end{column}
\end{columns}
\end{frame}

% ════════════════════════════════════════════════════════════
% SLIDE 5: Phase 3 — Co-folding
% ════════════════════════════════════════════════════════════
\begin{frame}{Phase 3: Co-folding with Boltz-2}
\vspace{-6pt}

\begin{columns}[T]
\begin{column}{0.55\textwidth}
\begin{block}{What Goes Into Each Run}
\begin{itemize}\small
\item \textbf{Protein sequence} (folds from scratch)
\item \textbf{Substrates} as SMILES\\
  {\scriptsize Adipic acid: \texttt{[O-]C(=O)CCCCC(=O)[O-]}}\\
  {\scriptsize HMDA: \texttt{[NH3+]CCCCCC[NH3+]}}
\item \textbf{ATP + Mg$^{2+}$} (CCD codes)
\item \textbf{Pocket constraints} from Phase 1--2
\end{itemize}
\end{block}

\vspace{2pt}
\colorbox{lightblue!50}{\parbox{0.92\textwidth}{\centering\small
\textbf{Cofactor normalization:} ANP/ACP $\rightarrow$ ATP, MN $\rightarrow$ MG\\
Model the pre-catalytic state, not crystallization artifacts
}}
\end{column}

\begin{column}{0.42\textwidth}
\begin{exampleblock}{Three Strategies}
\textbf{Cofactor-guided} {\small (28 synthetases)}
\begin{itemize}\scriptsize
\item ATP constrained to AlphaFill pocket
\item Substrates dock \textbf{freely} near cofactor
\item \textcolor{accentgreen}{Independent validation}
\end{itemize}

\vspace{2pt}
\textbf{Pocket-constrained} {\small (89 hydrolases)}
\begin{itemize}\scriptsize
\item Substrates anchored to predicted pocket
\item \textcolor{accentorange}{Circular with pocket selection}
\end{itemize}

\vspace{2pt}
\textbf{Blind} {\small (4 hydrolases)}
\begin{itemize}\scriptsize
\item No spatial constraints
\end{itemize}
\end{exampleblock}

\vspace{2pt}
\colorbox{lightgray}{\parbox{0.92\textwidth}{\centering\small
\textbf{122 complexes} co-folded\\(1 excluded: $\varepsilon$-PL synthetase, 1319 aa, GPU memory)
}}
\end{column}
\end{columns}
\end{frame}

% ════════════════════════════════════════════════════════════
% SLIDE 6: Phase 4 — Quality Metrics
% ════════════════════════════════════════════════════════════
\begin{frame}{Phase 4: How We Score Each Complex}
\vspace{-6pt}

\begin{columns}[T]
\begin{column}{0.48\textwidth}
\begin{block}{Binding Quality}
\begin{itemize}\scriptsize
\item \textbf{iPTM} --- overall interface confidence
\item \textbf{Substrate-specific iPTM} --- isolates protein--substrate confidence from ATP inflation
\item \textbf{SASA burial} --- fraction of substrate enclosed ($>$80\% = good)
\item \textbf{Steric clashes} --- severe ($<$1.5\AA), bad ($<$2.0\AA)
\end{itemize}
\end{block}

\begin{block}{Pocket Geometry}
\begin{itemize}\scriptsize
\item \textbf{Volume ratio} --- pocket vs.\ substrate (detect forced-fit or oversized)
\item \textbf{Forced-fit flags} --- groove binding, squeeze artifacts
\end{itemize}
\end{block}
\end{column}

\begin{column}{0.48\textwidth}
\begin{block}{Pocket Chemistry}
\begin{itemize}\scriptsize
\item \textbf{Holo pocket composition} --- \% hydrophobic, charged, polar within 5~\AA
\item \textbf{Catalytic Center Score} --- catalytic residue overlap with pocket
\end{itemize}
\end{block}

\begin{alertblock}{Apo-to-Holo Comparison}
\begin{itemize}\scriptsize
\item \textbf{Residue Jaccard} --- shared pocket-lining residues between AlphaFold (apo) and Boltz-2 (holo)
\item \textbf{Volume change} --- pocket expansion/contraction
\item High Jaccard = substrate at \textbf{predicted site}
\item Low Jaccard = substrate may be \textbf{elsewhere}
\end{itemize}
\end{alertblock}
\end{column}
\end{columns}

\vspace{2pt}
\colorbox{accentorange!15}{\parbox{0.92\textwidth}{\centering\scriptsize
\textbf{Why substrate-specific iPTM?} Overall iPTM = 0.913 (synthetase mean) vs.\ substrate iPTM = 0.893 --- ATP cofactor inflates the aggregate. Largest gap: PanC (0.962 vs.\ 0.914).
}}
\end{frame}

% ════════════════════════════════════════════════════════════
% SLIDE 7: Aggregate Results
% ════════════════════════════════════════════════════════════
\begin{frame}{Results: Synthetases vs.\ Hydrolases}
\vspace{-6pt}

\begin{columns}[T]
\begin{column}{0.40\textwidth}
\begin{center}
\scriptsize
\begin{tabular}{lcc}
\toprule
\textbf{Metric} & \textbf{Syn.} & \textbf{Hyd.} \\
& {\tiny ($n$=29)} & {\tiny ($n$=93)} \\
\midrule
Mean iPTM & \textbf{0.913} & 0.748 \\
Mean sub.\ iPTM & \textbf{0.893} & --- \\
Mean SASA (\%) & \textbf{86.1} & 61.5 \\
Catalytic ann.\ (\%) & \textbf{55} & 25 \\
Clash-free (\%) & 86 & \textbf{97} \\
\midrule
\multicolumn{3}{l}{\textit{Pocket geometry flags:}} \\
Oversized pocket & 39 & 7 \\
Pocket too small & 2 & \textcolor{accentred}{\textbf{41}} \\
Groove binding & 0 & \textcolor{accentred}{\textbf{23}} \\
\bottomrule
\end{tabular}
\end{center}
\end{column}

\begin{column}{0.57\textwidth}
\begin{exampleblock}{Why Synthetases Outperform}
\begin{itemize}\scriptsize
\item Selected for amide \textit{formation} $\rightarrow$ substrates fit naturally
\item Cofactor-guided co-folding provides strong constraints
\item Spacious pockets allow flexible substrate placement
\end{itemize}
\end{exampleblock}

\begin{alertblock}{Why Hydrolases Fail}
\begin{itemize}\scriptsize
\item Evolved for \textbf{charged} polymers (cyanophycin, polyglutamate)
\item \textbf{44\%} pockets too small for nylon monomers
\item \textbf{25\%} constrained to surface grooves
\item Poor metrics hold even with confirmed active-site constraints (65 candidates, mean iPTM = 0.821)
\end{itemize}
\end{alertblock}
\end{column}
\end{columns}
\end{frame}

% ════════════════════════════════════════════════════════════
% SLIDE 8: Top Synthetases Table
% ════════════════════════════════════════════════════════════
\begin{frame}{Top Synthetase Candidates}
\vspace{-6pt}
\begin{center}
\begin{adjustbox}{max width=\textwidth}
\small
\begin{tabular}{llccccc}
\toprule
\textbf{Gene} & \textbf{Family} & \textbf{iPTM} & \textbf{Sub.~iPTM} & \textbf{SASA (\%)} & \textbf{Hydrophobic (\%)} & \textbf{Charged (\%)} \\
\midrule
iucA & NIS synthetase & 0.990 & 0.986 & 89.6 & 30.1 & 26.5 \\
asbA & NIS synthetase & 0.990 & 0.983 & 98.7 & 20.7 & 36.2 \\
rhbC & NIS synthetase & 0.989 & 0.985 & 96.9 & 27.1 & 38.6 \\
\rowcolor{accentred!15} McbA & NIS synthetase & 0.987 & 0.979 & 93.4 & 42.9 & 25.5 \\
ddlB & ATP-grasp & 0.986 & 0.974 & 96.5 & 32.6 & 29.2 \\
sbnE & NIS synthetase & 0.985 & 0.975 & 87.4 & 28.6 & 37.7 \\
pylC & ATP-grasp & 0.983 & 0.968 & 92.7 & 37.0 & 30.0 \\
murC & Mur ligase & 0.981 & 0.966 & 70.5 & 25.0 & 32.7 \\
ddaF & ANL superfamily & 0.980 & 0.963 & 78.6 & 39.7 & 25.0 \\
murD & Mur ligase & 0.975 & 0.957 & 80.3 & 31.6 & 26.3 \\
\bottomrule
\end{tabular}
\end{adjustbox}
\end{center}

\vspace{4pt}
\begin{columns}[T]
\begin{column}{0.48\textwidth}
\colorbox{lightblue!50}{\parbox{0.92\textwidth}{\centering\small
\textbf{NIS synthetases} rank highest --- native function is mechanistically identical to target
}}
\end{column}
\begin{column}{0.48\textwidth}
\colorbox{accentred!15}{\parbox{0.92\textwidth}{\centering\small
\textbf{McbA} ranks 4th but is \textbf{experimentally inactive} --- iPTM alone doesn't predict activity
}}
\end{column}
\end{columns}
\end{frame}

% ════════════════════════════════════════════════════════════
% SLIDE 9: Substrate Compatibility
% ════════════════════════════════════════════════════════════
\begin{frame}{Substrate Compatibility: Chemistry Meets Structure}
\vspace{-8pt}
\begin{center}
\includegraphics[width=0.50\textwidth]{analysis_results/substrate_compatibility.png}
\end{center}
\vspace{-6pt}
\begin{center}\scriptsize
Native substrate chain length vs.\ holo pocket hydrophobicity.
\textbf{Green zone}: long-chain + hydrophobic pocket = best for aliphatic nylon monomers.
\end{center}
\end{frame}

% ════════════════════════════════════════════════════════════
% SLIDE 10: Pocket Chemistry 2×2
% ════════════════════════════════════════════════════════════
\begin{frame}{Pocket Chemistry: Apo vs.\ Holo}
\vspace{-8pt}
\begin{center}
\includegraphics[width=0.52\textwidth]{analysis_results/pocket_chemistry_landscape.png}
\end{center}
\vspace{-6pt}
\begin{columns}[T]
\begin{column}{0.48\textwidth}
\centering\scriptsize
\textbf{Synthetases (A, B):} PanC at highest hydrophobicity in both apo and holo.
\end{column}
\begin{column}{0.48\textwidth}
\centering\scriptsize
\textbf{Hydrolases (C, D):} NylC ({\color{accentgreen}$\bigstar$}) clear outlier --- 42\% hydrophobic, 5\% charged vs.\ mean 31\%/29\%.
\end{column}
\end{columns}
\end{frame}

% ════════════════════════════════════════════════════════════
% SLIDE 11: Apo-to-Holo
% ════════════════════════════════════════════════════════════
\begin{frame}{Apo vs.\ Holo: Does the Substrate Reach the Catalytic Site?}
\vspace{-8pt}

\begin{center}
\small
\begin{tabular}{lcccl}
\toprule
\textbf{Candidate} & \textbf{Jaccard} & \textbf{Vol.\ $\Delta$} & \textbf{Sub.\ iPTM} & \textbf{Interpretation} \\
\midrule
\rowcolor{accentgreen!15} PanC & \textbf{0.889} & +0.1\% & 0.914 & Substrate in catalytic pocket; no rearrangement \\
\rowcolor{accentgreen!15} BacD & \textbf{0.914} & +3.5\% & 0.907 & Stable pocket; high confidence \\
\rowcolor{lightblue!30} gss & 0.812 & +5.3\% & 0.936 & Moderate stability \\
\rowcolor{lightblue!30} ddaF & 0.829 & $-$7.7\% & 0.963 & Moderate stability; slight contraction \\
\rowcolor{accentred!15} McbA & 0.883 & $-$19.4\% & 0.979 & Correct site \textit{but inactive} (geometry mismatch) \\
\rowcolor{accentorange!15} GshA & \textbf{0.025} & +998\% & 0.959 & \textbf{Completely different site} \\
\bottomrule
\end{tabular}
\end{center}

\vspace{4pt}
\begin{columns}[T]
\begin{column}{0.48\textwidth}
\begin{exampleblock}{High Jaccard = Independent Validation}
\begin{itemize}\small
\item Apo pocket (AlphaFill) $\approx$ holo pocket (Boltz-2)
\item Two independent methods \textbf{agree} on binding site
\item PanC and BacD: substrate at the predicted catalytic pocket without conformational reorganization
\end{itemize}
\end{exampleblock}
\end{column}
\begin{column}{0.48\textwidth}
\begin{alertblock}{GshA: A Cautionary Case}
\begin{itemize}\small
\item Jaccard = 0.025: almost \textbf{zero overlap}
\item Boltz-2 placed substrate at a \textbf{different site}
\item Still ranked \#2 on substrate chemistry\\(glutamate $\approx$ adipic acid)
\item Experiment will test: does substrate analogy alone predict activity?
\end{itemize}
\end{alertblock}
\end{column}
\end{columns}
\end{frame}

% ════════════════════════════════════════════════════════════
% SLIDE 12: McbA Case Study
% ════════════════════════════════════════════════════════════
\begin{frame}{McbA: Why Good Metrics $\neq$ Activity}
\vspace{-6pt}

\begin{columns}[T]
\begin{column}{0.48\textwidth}
\begin{block}{McbA's Metrics Look Great}
\begin{itemize}\scriptsize
\item iPTM = 0.987 (\textbf{4th highest})
\item SASA = 93.4\%, Hydrophobicity = 42.9\%
\item Jaccard = 0.883 (correct pocket)
\item Catalytic Center Score = 100\%
\end{itemize}
\end{block}

\colorbox{accentred!20}{\parbox{0.92\textwidth}{\centering
\normalsize\textbf{Result: No activity}\\[1pt]
\scriptsize Expressed in \textit{E. coli}, no amide synthesis detected
}}

\vspace{4pt}
\begin{exampleblock}{The Lesson}
\begin{itemize}\scriptsize
\item \textbf{What} is in the pocket $\neq$ \textbf{where} residues are
\item Need \textbf{catalytic geometry}: near-attack distances \& angles
\item $\rightarrow$ Key motivation for proposed work
\end{itemize}
\end{exampleblock}
\end{column}

\begin{column}{0.48\textwidth}
\begin{alertblock}{Structural Explanation}
\begin{itemize}\scriptsize
\item McbA's native substrates: \textbf{aromatic/branched} carboxylic acids (microcystin pathway)
\item Target substrates: \textbf{linear aliphatic} (adipic acid, HMDA)
\item Pocket \textit{composition} is similar but \textit{shape} is different
\item Cannot orient adipic acid's linear chain for adenylation
\end{itemize}
\end{alertblock}

\vspace{4pt}
\colorbox{lightblue!50}{\parbox{0.92\textwidth}{\centering\scriptsize
\textbf{Pocket composition} (42.9\% hydrophobic) $\checkmark$\\
\textbf{Substrate placement} (Jaccard = 0.883) $\checkmark$\\
\textbf{Catalytic geometry} (substrate orientation) \textcolor{accentred}{\textbf{$\times$}}\\[2pt]
$\rightarrow$ All current metrics pass; the missing descriptor is \textbf{spatial arrangement} of catalytic residues relative to substrate reactive centers
}}
\end{column}
\end{columns}
\end{frame}

% ════════════════════════════════════════════════════════════
% SLIDE 13: NylC Benchmark + Hydrolase Control
% ════════════════════════════════════════════════════════════
\begin{frame}{Hydrolase Control: Why 93 Screened Hydrolases Failed}
\vspace{-4pt}
\begin{columns}[T]
\begin{column}{0.48\textwidth}
\begin{block}{Families Screened (None Active)}
\begin{itemize}\small
\item 24 poly-$\gamma$-glutamate hydrolases
\item 22 cyanophycinases
\item 7 pseudomurein endoisopeptidases
\item 5 zinc-dependent hydrolases
\item \textbf{None} use NylC's Ntn-hydrolase mechanism
\end{itemize}
\end{block}

\begin{alertblock}{Two Independent Failures}
\begin{enumerate}\small
\item \textbf{Pocket chemistry:} 29\% charged (vs.\ NylC's 5\%) --- evolved for charged polymers, not aliphatic nylon
\item \textbf{Pocket geometry:} 44\% too small + 25\% groove-bound --- fundamentally cannot accommodate nylon monomers
\end{enumerate}
\end{alertblock}
\end{column}

\begin{column}{0.48\textwidth}
\begin{exampleblock}{NylC: The Benchmark}
\begin{itemize}\small
\item Only enzyme with confirmed nylon-6,6 endo-hydrolysis
\item Co-folded in \textbf{autocleavage form} (split at Thr267)
\item Holo: 42.1\% hydrophobic, 5.3\% charged
\item Clear outlier in pocket chemistry landscape
\end{itemize}
\end{exampleblock}

\vspace{4pt}
\colorbox{lightblue!50}{\parbox{0.92\textwidth}{\centering\small
\textbf{Takeaway:} The 93 hydrolases serve as a \textbf{structural negative control} --- their poor metrics across \textit{both} chemistry and geometry confirm that nylon compatibility requires a specific pocket environment, not just amide-cleaving activity
}}
\end{column}
\end{columns}
\end{frame}

% ════════════════════════════════════════════════════════════
% SLIDE 14: Recommended Test Set
% ════════════════════════════════════════════════════════════
\begin{frame}{Recommended Experimental Test Set}
\vspace{-6pt}

\begin{center}
\begin{adjustbox}{max width=\textwidth}
\small
\begin{tabular}{clccccl}
\toprule
\textbf{\#} & \textbf{Gene} & \textbf{Sub.\ iPTM} & \textbf{SASA} & \textbf{Hydro.} & \textbf{Jaccard} & \textbf{Why this candidate?} \\
\midrule
1 & \textbf{PanC} & 0.914 & 97.8\% & \textbf{48.2\%} & \textbf{0.889} & Pantoate $\approx$ adipic acid; 100\% cat.\ center; \textit{E. coli} \\
2 & \textbf{GshA} & 0.959 & 82.4\% & 29.0\% & \textcolor{accentred}{0.025} & Glutamate (147 Da) $\approx$ adipic acid (146 Da); \textit{E. coli} \\
3 & \textbf{BacD} & 0.907 & 93.5\% & 40.7\% & \textbf{0.914} & Proven broad specificity; 8 mutagenesis variants \\
4 & \textbf{gss} & 0.936 & 95.3\% & 33.8\% & 0.812 & Phosphorylation mechanism (diversity); \textit{E. coli} \\
5 & \textbf{DdaF} & 0.963 & 78.6\% & 39.7\% & 0.829 & Highest sub.\ iPTM; follow-up candidate \\
\bottomrule
\end{tabular}
\end{adjustbox}
\end{center}

\vspace{4pt}
\begin{columns}[T]
\begin{column}{0.32\textwidth}
\begin{exampleblock}{\small Strongest Structural}
\centering\small
\textbf{PanC}: highest hydrophobicity, most stable pocket, 100\% catalytic overlap
\end{exampleblock}
\end{column}
\begin{column}{0.32\textwidth}
\begin{alertblock}{\small Strongest Chemical}
\centering\small
\textbf{GshA}: best MW match but low Jaccard --- will test if analogy alone predicts activity
\end{alertblock}
\end{column}
\begin{column}{0.32\textwidth}
\begin{block}{\small Best Engineered}
\centering\small
\textbf{BacD}: known mutations + most stable pocket (Jaccard = 0.914)
\end{block}
\end{column}
\end{columns}

\vspace{4pt}
\colorbox{lightgray}{\parbox{0.94\textwidth}{\centering\small
\textbf{Previously tested:} McbA (inactive --- geometry mismatch) + PylC (failed expression --- archaeal protein)\\
$\rightarrow$ Combined: \textbf{7 candidates} across \textbf{5 mechanistic families}
}}
\end{frame}

% ════════════════════════════════════════════════════════════
% SLIDE 15: Next Steps
% ════════════════════════════════════════════════════════════
\begin{frame}{Next Steps}
\vspace{-4pt}
\begin{columns}[T]
\begin{column}{0.48\textwidth}
\begin{block}{Experimental (Aim 3)}
\begin{itemize}
\item Express \& assay \textbf{5 candidates} in \textit{E. coli}
\item Each tests a \textbf{structural hypothesis}:
  \begin{itemize}\small
  \item PanC: substrate similarity + pocket stability
  \item GshA: MW analogy vs.\ low Jaccard
  \item BacD: broad specificity + stable pocket
  \item gss: alternative activation mechanism
  \item DdaF: high iPTM, weak chemistry
  \end{itemize}
\item Feed results back to agent
\end{itemize}
\end{block}
\end{column}

\begin{column}{0.48\textwidth}
\begin{exampleblock}{Computational (Aim 2)}
\begin{itemize}
\item \textbf{Catalytic geometry:} near-attack distances \& B\"urgi--Dunitz angles from co-folded structures
\item \textbf{Spatially resolved descriptors:} per-residue distance from substrate reactive centers
\item \textbf{Cross-family comparison:} how do NIS, ATP-grasp, and Mur ligase folds position the same substrates?
\end{itemize}
\end{exampleblock}
\end{column}
\end{columns}

\vspace{6pt}
\colorbox{accentgreen!15}{\parbox{0.88\textwidth}{\centering
\textbf{The question is not which candidates are active, but which structural hypotheses hold} --- revealing whether substrate similarity, pocket composition, or catalytic geometry best predicts activity on non-natural substrates.
}}
\end{frame}

% ════════════════════════════════════════════════════════════
% SLIDE 16: Summary
% ════════════════════════════════════════════════════════════
{
\setbeamercolor{background canvas}{bg=deepblue}
\begin{frame}[plain]
\vfill
\begin{center}
{\Large\bfseries\color{white} Summary}\\[16pt]
{\normalsize\color{lightblue}
\begin{tabular}{rl}
\color{accentorange}{122} & \color{white} enzyme--substrate complexes co-folded (29 synthetases + 93 hydrolases) \\[6pt]
\color{accentorange}{8} & \color{white} independent quality metrics per complex (no composite score) \\[6pt]
\color{accentorange}{3$\rightarrow$22} & \color{white} hydrolase catalytic annotations via FoldSeek + LLM \\[6pt]
\color{accentorange}{44\%} & \color{white} of hydrolases have pocket geometry incompatible with nylon \\[6pt]
\color{accentorange}{5} & \color{white} candidates prioritized for experimental testing \\[6pt]
\color{accentorange}{1} & \color{white} key lesson: McbA shows that binding $\neq$ catalysis \\[6pt]
\end{tabular}
}
\end{center}
\vfill
\end{frame}
}

\end{document}
